%% LyX 2.5.1 created this file.  For more info, see https://www.lyx.org/.
%% Do not edit unless you really know what you are doing.
\documentclass[english,hebrew]{article}
\usepackage[HE8,T1]{fontenc}
\usepackage[utf8]{inputenc}
\usepackage{amsmath}
\usepackage{amssymb}
\usepackage{babel}
\begin{document}
\title{מבוא לטופולוגיה קבוצתית: קבוצות פתוחות במרחבים מטריים}
\maketitle
\begin{description}
\item [{קטגוריות:}] טופולוגיה
\item [{תגים:}] מרחבים מטריים, מרחבים טופולוגיים
\item [{מזהה:}] \L{open\_sets\_in\_metric\_spaces}
\end{description}
בפוסט הקודם ראינו מה זה \textbf{מרחב מטרי}. מה שטוב במרחב מטרי זה
שהוא מהווה הכללה די סבירה של מה שאנחנו עושים בחשבון אינפיניטסימלי,
ומהווה קרש קפיצה טוב אל ההכללה הבאה. בפוסט הקודם השתמשתי במושג הגבול
של \textbf{סדרה} כדי לתת מוטיבציה להכללה, הפעם בואו נלך עוד צעד קדימה
אל מושג הגבול של \textbf{פונקציה}.

ההגדרה הולכת כך: אם \L{$f:\mathbb{R}\to\mathbb{R}$} היא פונקציה כלשהי,
אנחנו כותבים \L{$\lim_{x\to x_{0}}f\left(x\right)=L$} )\textquotedblleft הפונקציה
\L{$f$} שואפת אל \L{$L$} כאשר \L{$x$} שואף אל \L{$x_{0}$}( כדי
לתאר את המצב שבו \textbf{לכל} \L{$\varepsilon>0$} \textbf{קיים} \L{$\delta>0$}
כך שאם \L{$0<\left|x-x_{0}\right|<\delta$} אז \L{$\left|f\left(x\right)-L\right|<\varepsilon$}.

זו הגדרה קשה לעיכול במבט ראשון אבל אני מניח שאנחנו כבר מכירים אותה
)שימו לב לדקות הזו של \L{$0<\left|x-x_{0}\right|$} שבאה לציין במפורש
שאפשר להניח ש-\L{$x\ne x_{0}$}; בלי הדרישה הזו יש לנו את ההגדרה של
\textbf{רציפות} ב-\L{$x_{0}$}(. השאלה היא איך להכליל אותה בצורה נוחה.
ראשית, הנה דרך אחרת לכתוב את אותה הגדרה:

\textbf{לכל} \L{$\varepsilon>0$} \textbf{קיים} \L{$\delta>0$} כך
שאם \L{$x\in\left(x_{0}-\delta,x_{0}+\delta\right)$} ו-\L{$x\ne x_{0}$}
אז \L{$f\left(x\right)\in\left(L-\varepsilon,L+\varepsilon\right)$}.

הכנסנו כאן לתמונה משהו חדש - \textbf{קטע פתוח}. באופן כללי קטע פתוח
הוא הקבוצה \L{$\left(a,b\right)=\left\{ x\in\mathbb{R}\ |\ a<x<b\right\} $}.
עבור גבול באופן ספציפי הקטע הפתוח הוגדר בצורה סימטרית שכזו ביחס לנקודת
אמצע - זה כבר מזכיר את המושג מהפוסט הקודם של \textquotedblleft כדור
פתוח\textquotedblright . בואו נזכיר אותו. בפוסט הקודם דיברנו על מרחב
מטרי כללי \L{$\left(X,d\right)$} ואז בהינתן נקודה \L{$a\in X$} ומספר
ממשי אי שלילי \L{$r\ge0$}, הכדור הפתוח מרדיוס \L{$r$} סביב \L{$a$}
היה

\L{$B_{d}\left(a,r\right)=\left\{ x\in X\ |\ d\left(a,x\right)<r\right\} $}

זה מאפשר לנסח מחדש את הגדרת הגבול של פונקציה באופן כללי יותר. בואו
נניח שיש לנו \textbf{שני} מרחבים מטריים, \L{$\left(X,d_{X}\right),\left(Y,d_{Y}\right)$}
ופונקציה \L{$f:X\to Y$}. אנחנו מסתכלים על נקודה \L{$x_{0}\in X$}
וערך \L{$L\in Y$} ומסמנים \L{$\lim_{x\to x_{0}}f\left(x\right)=L$}
אם \textbf{לכל} \L{$\varepsilon>0$} \textbf{קיים} \L{$\delta>0$}
כך שאם \L{$x\in B_{d_{X}}\left(x_{0},\delta\right)$} ו-\L{$x\ne x_{0}$}
אז \L{$f\left(x\right)\in B_{d_{Y}}\left(L,\varepsilon\right)$}.

יפה, התקדמנו. אבל רגע, בואו ננסה להיזכר מה קורה בחשבון אינפיניטסימלי
- איך אנחנו משתמשים בהגדרת הגבול? האם זה קריטי לנו לקחת דווקא \textbf{כדור}
סביב נקודה? ובכן, לא. לא באמת צריך לקחת קבוצה עם מבנה \textquotedblleft יפה\textquotedblright .
יותר מזה - הקטע של לומר \textquotedblleft לכל \L{$\varepsilon>0$}\textquotedblright{}
הוא קצת תובעני מדי - לא באמת חשוב לנו שההגדרה תעבוד \textbf{לכל} אפסילון,
פשוט לאפסילון שהוא \textquotedblleft קטן כרצוננו\textquotedblright .
כלומר, לא משנה כמה קטן ניקח את \L{$\varepsilon$}, עדיין יהיה \L{$\delta$}
שעובד עבורו. מן הסתם אם \L{$\delta$} עובד עבור \L{$\varepsilon$}
\textbf{מסוים} אז הוא עובד עבור \textbf{כל} אפסילון אחר שגדול או שווה
אליו; מה שאנחנו באמת צריכים הוא פשוט סדרה של אפסילונים שמתקרבים לאפס
\textquotedblleft כרצוננו\textquotedblright .

זה מוביל אותנו אל מילה אמורפית יותר, \textquotedblleft סביבה\textquotedblright .
כשלמדתי אינפי בפעם הראשונה, קבוצה כמו \L{$\left(x_{0}-\delta,x_{0}+\delta\right)$}
נקראה \textquotedblleft סביבת-\L{$\delta$} של \L{$x_{0}$}\textquotedblright ,
אבל אפשר היה לדבר על סביבה של נקודה גם בלי לקחת לה רדיוס ספציפי. בלי
המבנה הספציפי הזה של רדיוס ומרחק, מה \textquotedblleft סביבה\textquotedblright{}
של \L{$x_{0}$} עדיין צריכה לקיים? היא מן הסתם צריכה להכיל את \L{$x_{0}$}
עצמו בתור איבר )אם כי באינפי אהבו גם לדבר על \textquotedblleft סביבה
\textbf{נקובה}\textquotedblright{} של \L{$x_{0}$} כדי להסתכל רק על
\L{$x$}-ים בסביבה של \L{$x_{0}$} שלא שווים לו, אבל נעזוב את זה -
בהגדרות שלי בכוונה הלכתי על \textquotedblleft איבר \L{$x$} בסביבה
של \L{$x_{0}$} כך ש-\L{$x\ne x_{0}$}\textquotedblright (, אבל מה
עוד? ה\textquotedblright עוד\textquotedblright{} הזה נקרא \textbf{קבוצה
פתוחה} ובגלל שזה הולך להפוך למושג המרכזי שלנו, בואו ננסה להבין מאיפה
זה בכלל מגיע.

אם כן, ראשית בואו ננסה לתת הגדרה חדשה לגבול עם המושג הזה של \textquotedblleft סביבה\textquotedblright :

נתונה פונקציה \L{$f:X\to Y$}. אנחנו מסתכלים על נקודה \L{$x_{0}\in X$}
וערך \L{$L\in Y$} ומסמנים \L{$\lim_{x\to x_{0}}f\left(x\right)=L$}
אם \textbf{לכל} סביבה \L{$B\subseteq Y$} של \L{$L$} \textbf{קיימת}
סביבה \L{$A\subseteq X$} של \L{$x_{0}$}, כך שאם \L{$x\in A$} ו-\L{$x\ne x_{0}$}
אז \L{$f\left(x\right)\in B$}.

אם כל הדרישה שלנו מ-\L{$A,B$} הוא שיתקיים \L{$x_{0}\in A,L\in B$}
וזהו, ההגדרה הופכת להיות \textbf{טריוויאלית}. בואו ניקח \L{$f$} כלשהי,
לא משנה כמה מוזרה, וניקח \L{$x_{0},L$} כלשהו. אין שום סיבה שיתקיים
\L{$\lim_{x\to x_{0}}f\left(x\right)=L$} במקרה הכללי הזה. אבל כעת,
בואו ניקח סביבה \L{$B$} כלשהי של \L{$L$} ועבורה נבחר את הסביבה הספציפית
\L{$A=\left\{ x_{0}\right\} $}. כלומר, פשוט אין בסביבה הזו שום דבר
חוץ מ-\L{$x_{0}$} עצמה. אין בה \L{$x\in A$} כך ש-\L{$x\ne x_{0}$}
ולכן כל התנאי של \textquotedblleft אם \L{$x\in A$} ו-\L{$x\ne x_{0}$}
אז \L{$f\left(x\right)\in B$}\textquotedblright{} מתקיים \textbf{באופן
ריק}. זה בלתי אפשרי. לב העניין פה הוא שסביבה של \L{$x_{0}$} תהיה
חייבת להיות \textquotedblleft לא טריוויאלית\textquotedblright ; לקבץ
אליה איברים קרובים ל-\L{$x_{0}$} אם יש כאלו.

שימו לב ל\textquotedblright אם יש כאלו\textquotedblright . לא תמיד
חייבים להיות. יש דוגמא קלאסית למטריקה שנקראת \textbf{המטריקה הדיסקרטית}
שבה פשוט מגדירים, נאמר, \L{$d\left(a,b\right)=\begin{cases}
1 & a\ne b\\
0 & a=b
\end{cases}$}. במקרה הזה, \L{$B\left(a,\frac{1}{2}\right)=\left\{ a\right\} $}
וכדורים פתוחים כאמור \textbf{כן} מתאימים למה שזה לא יהיה שסביבה אמורה
להיות )כי עם כדורים פתוחים כבר קיבלנו הגדרה שקולה להגדרת הגבול המקורית(.
כלומר אנחנו יכולים בעיקרון לקבל סביבות של \L{$a$} שכוללות רק את \L{$a$}
עצמו, אם כי זה אומר משהו די מהותי על האופי של המרחב - אבל זה שונה
מאשר סתם לקחת את \L{$a$} וזהו. שימו לב שעם כדורים פתוחים אנחנו לא
מסוגלים לקבל את הקבוצה הזו באופן כללי: אפילו אם אסתכל על \L{$B\left(a,0\right)=\left\{ x\in X\ |\ d\left(a,x\right)<0\right\} $},
מה שאני אקבל פה הוא את הקבוצה הריקה. כדי לקבל את הקבוצה שכוללת רק
את \L{$a$} אני צריך להפוך את ה-\L{$<$} שבפנים אל \L{$\le$}, מה
שהופך את מה שאנחנו קוראים לו \textquotedblleft כדור \textbf{פתוח}\textquotedblright{}
אל מה שנקרא \textquotedblleft כדור \textbf{סגור}\textquotedblright{}
וזה אכן כבר סיפור שונה שעוד מעט נדבר גם עליו.

אוקיי, אז אנחנו רוצים להגדיר איכשהו קבוצה \L{$A$} שמקיימת את התכונה
שאם \L{$a\in A$}, אז יש בה \textquotedblleft איברים קרובים ל-\L{$a$}\textquotedblright ,
אבל גם את זה צריך לחדד - כדי שהגדרת הגבול תעבוד, לא מספיק לקחת שניים-שלושה
איברים שקרובים ל-\L{$a$} ולהוסיף אותם. אנחנו נצטרך את \textbf{כל}
האיברים שקרובים ל-\L{$a$} החל מרמת קרבה מסוימת. וזה... בדיוק מה שכדור
פתוח עשה. אז למה שלא נשתמש בו?

זה מוביל אותנו להגדרה הבאה: קבוצה \L{$A$} היא \textbf{פתוחה} אם לכל
\L{$a\in A$} קיים \L{$r>0$} כך ש-\L{$B\left(a,r\right)\subseteq A$},
כלומר לכל נקודה בקבוצה יש כדור פתוח לא ריק סביבה שמוכל כולו בתוך הקבוצה.

האם יש לנו דוגמאות קונקרטיות לקבוצה פתוחה? בוודאי - כדור פתוח הוא
בעצמו קבוצה פתוחה, ומכיוון שנשתמש בזה בהמשך בואו נוכיח את זה: אם \L{$B\left(a,r\right)$}
הוא כדור פתוח כלשהו ו-\L{$x\in B\left(a,r\right)$} היא נקודה כלשהי
בכדור הפתוח, אז \L{$d\left(a,x\right)<r$}. בואו ניקח כדור פתוח סביב
\L{$x$} שכולו מוכל ב-\L{$B\left(a,r\right)$}, כלומר אף נקודה בו
לא רחוקה מ-\L{$a$} יותר מאשר \L{$r$}. כמה גדול צריך הכדור הפתוח
הזה להיות? בואו נסמן את הגודל שלו ב-\L{$s$} לבינתיים. עכשיו ניקח
נקודה \L{$y\in B\left(x,s\right)$} ונשתמש באי-שוויון המשולש כדי לחסום
את המרחק שלה מ-\L{$a$}:

\L{$d\left(y,a\right)\le d\left(y,x\right)+d\left(x,a\right)<s+d\left(x,a\right)$}.
אז כדי לקבל \L{$d\left(y,a\right)<r$} צריך שיתקיים \L{$s+d\left(x,a\right)<r$},
כלומר \L{$s<r-d\left(x,a\right)$} וזה אפשרי לבחור \L{$s$} כזה כי
\L{$d\left(a,x\right)<r$} ולכן \L{$0<r-d\left(a,x\right)$}.

עכשיו אפשר להגדיר \textbf{סביבה} של \L{$a$} בתור קבוצה פתוחה שמכילה
את \L{$a$}. הגדרת הגבול שנתתי למעלה עם \textquotedblleft סביבה\textquotedblright{}
עובדת עכשיו כמו שצריך )זה תרגיל טוב להוכיח שהיא שקולה למה שהראינו
קודם( ו... נראה שסתם סיבכנו? קודם ההגדרה שלי התבססה על כדורים פתוחים,
ועכשיו היא מסתמכת על קבוצות פתוחות שבתורן מסתמכות על כדורים? זה לא
כזה שיפור.

אני מסכים, כרגע זה לא נראה עדיין כמו שיפור. אז בואו נשחק קצת עם המושג
החדש שהמצאנו ונראה איזה דברים מתקשרים אליו.

ראשית, דיברתי על גבול של פונקציה אבל אפשר לדבר גם על גבול של סדרה.
אם \L{$\left(X,d\right)$} הוא מרחב מטרי כלשהו ו-\L{$a_{1},a_{2},a_{3},\ldots$}
היא סדרת איברים של \L{$X$}, אז נגיד ש-\L{$\lim_{n\to\infty}a_{n}=a$}
עבור \L{$a\in X$} אם לכל סביבה של \L{$a$} )כלומר - כל קבוצה פתוחה
ש-\L{$a$} שייך אליה(, קיים \L{$N$} כך שכל אברי הסדרה החל מ-\L{$a_{N}$}
שייכים לסביבה. אני חייב להודות שאני אוהב את הניסוח הזה - כבר אין צורך
לדבר במפורש על מרחקים ואפסילונים ואפילו לא על כדורים. כל יותר... אה...
נקי? זה לא בהכרח טוב יותר - הגישה האבסטרקטית הזו עשויה לפעמים להקשות
על יצירת אינטואיציה לגבי \textquotedblleft מה הולך פה ואיך נוכיח את
הדבר הזה והזה\textquotedblright{} אבל זה נותן נקודת מבט שמרגישה \textbf{שונה}
וזה חשוב.

אחד מהדברים הראשונים שרואים באינפי )אם מתחילים עם גבולות של סדרות(
זה שאם לסדרה קיים גבול אז הוא \textbf{יחיד}. האם זה נכון גם כאן? נניח
שלסדרה \L{$a_{1},a_{2},a_{3},\ldots$} יש שני גבולות שנסמן ב-\L{$a,b$}
וננסה להוכיח ש-\L{$a=b$}. 

באינפי ההוכחה מאוד חשבונית באופי שלה: לוקחים \L{$\varepsilon>0$}
כלשהו, מוצאים בעזרת הגדרת הגבול איבר \L{$x$} בסדרה שמקיים \L{$\left|a-x\right|<\frac{\varepsilon}{2},\left|b-x\right|<\frac{\varepsilon}{2}$}
)כלומר - אנחנו משתמשים בהגדרת הגבול כשאנחנו \textquotedblleft מזינים
לתוכה\textquotedblright{} את \L{$\frac{\varepsilon}{2}$} פעם ביחס
ל-\L{$a$} ופעם ביחס ל-\L{$b$}( ואז משתמשים באי-שוויון המשולש כדי
לקבל 

\L{$\left|a-b\right|\le\left|a-x\right|+\left|x-b\right|<\frac{\varepsilon}{2}+\frac{\varepsilon}{2}=\varepsilon$}

בצורה הזו קיבלנו ש-\L{$\left|a-b\right|$} קטן \textbf{מכל} מספר גדול
מאפס, לכן \L{$\left|a-b\right|=0$} ומכאן \L{$a=b$}. איך בכלל מתחילים
להעביר את כל זה לעולם המונחים של \textquotedblleft סביבות\textquotedblright ?
ובכן, לא בלי הנחות נוספות. חלק מהעניין עם אבסטרקציות הוא שהן מסוגלות
לתפוס יותר מאשר את מה שעשינו לו אבסטרקציה, כך שטענות שהיו נכונות קודם
יפסיקו לעבוד לגמרי. יחידות הגבול היא אחת מהטענות הללו. בואו נראה מה
בעצם המהות של ההוכחה באינפי - אנחנו אומרים \textquotedblleft בואו
ניקח סביבות קטנות ממש של \L{$a,b$}. בגלל שהסדרה מתכנסת לשני הגבולות
הללו, יהיה איבר שנמצא בשתי הסביבות הללו ואז - בום!\textquotedblright .
אם כן, זה לב העניין פה. אפשר לחלק את ההוכחה לשני חלקים:
\begin{enumerate}
\item מוכיחים שבמרחב \L{$X$} שלנו, לכל שתי נקודות \L{$a\ne b$} קיימות
שתי סביבות \L{$a\in A,b\in B$} כך ש-\L{$A\cap B=\emptyset$}.
\item מוכיחים שאם אותה סדרה מתכנסת גם ל-\L{$a$} וגם ל-\L{$b$} עבור \L{$a\ne b$},
אז קיימות שתי סביבות \L{$a\in A,b\in B$} כך ש-\L{$A\cap B\ne\emptyset$}
)בסתירה ל-{\beginL 1\endL}, מה שמפיל את ההנחה שאותה סדרה מתכנסת לשתי
נקודות שונות(.
\end{enumerate}
התכונה בשלב {\beginL 1\endL}, \textquotedblleft כל שתי נקודות שונות
ניתנות להפרדה על ידי סביבות שונות\textquotedblright{} היא כל כך חשוב
שיש לה סימון, \L{T2}, ויש לה שם - מרחב \L{$X$} שמקיים אותה נקרא \textbf{מרחב
האוסדורף}. בשלב הזה כל מה שאנחנו מכירים הוא מרחבים מטריים ועוד שניה
אני אוכיח ש\textbf{כל} המרחבים המטריים הם האוסדורפיים, ולכן לא לגמרי
ברור איזה מרחב יכול \textbf{לא} לקיים אותה; על זה נדבר בהמשך.

עכשיו, מה שקורה בשלב {\beginL 2\endL} בעצם אומר לנו \textquotedblleft במרחב
האוסדורף, לסדרה יש גבול יחיד\textquotedblright{} ואת זה \textbf{כן}
אפשר להוכיח רק עם הטרמינולוגיה של סביבות: ההאוסדורפיות של המרחב מספרת
לנו שיש שתי סביבות \L{$a\in A,b\in B$} כך ש-\L{$A\cap B=\emptyset$}.
הגדרת ההתכנסות מספרת לנו שקיים \L{$N_{A}$} כך שלכל \L{$n\ge N_{A}$},
מתקיים ש-\L{$a_{n}\in A$}. בדומה קיים \L{$N_{B}$}. ניקח \L{$n\ge\max\left\{ N_{A}.N_{B}\right\} $}
והוא יקיים \L{$a_{n}\in A$} וגם \L{$a_{n}\in B$}, כך ש-\L{$A\cap B\ne\emptyset$}
וזו הסתירה שרצינו. הוכחה \textquotedblleft נקיה\textquotedblright{}
בלי חשבונות קטנוניים; לאן זה הלך?

זה הלך, כמובן, להוכחה של שלב {\beginL 1\endL} - להוכחה שמרחב מטרי
הוא האוסדורפי. אנחנו מניחים ש-\L{$a\ne b$}, אז \L{$d\left(a,b\right)>0$}
)כי אחת מהאקסיומות של מרחבים מטריים היא ש-\L{$d\left(a,b\right)=0\iff a=b$}(.
נסמן \L{$\varepsilon=d\left(a,b\right)$}. עכשיו, נגדיר \L{$A=B\left(a,\frac{\varepsilon}{2}\right),B=B\left(a,\frac{\varepsilon}{2}\right)$}
)אוי לא... תסלחי לי על השימוש בכפול ב-\L{$B$} גם לתיאור סביבה קונקרטית
וגם לתיאור \textquotedblleft כדור פתוח\textquotedblright (. עכשיו,
ראינו שכדורים פתוחים הם קבוצות פתוחות ולכן באמת קיבלנו סביבות, אבל
למה הן זרות? נניח ש-\L{$c\in A\cap B$} אז \L{$d\left(a,c\right)<\frac{\varepsilon}{2}$}
וגם \L{$d\left(b,c\right)<\frac{\varepsilon}{2}$} ולכן לפי אי שוויון
המשולש, 

\L{$\varepsilon=d\left(a,b\right)\le d\left(a,c\right)+d\left(b,c\right)<\frac{\varepsilon}{2}+\frac{\varepsilon}{2}=\varepsilon$}

וקיבלנו ש-\L{$\varepsilon<\varepsilon$} - סתירה.

זה מסיים את ההוכחה, ועכשיו \textbf{שוב} נראה שלא עשינו כלום! כל המריבים
של ההוכחה המקורית עדיין כאן, ובנוסף יש עוד מושגים שזרקנו פנימה! אבל
לדעתי מה שעשינו פה היה שונה מהותית, בגלל הפיצול בין שלבים {\beginL 1\endL}
ו-{\beginL 2\endL}. שלב {\beginL 1\endL} נותן לנו תכונה כללית שיכולה
להיות במרחבים, ושלב {\beginL 2\endL} הוא לב לבה של הטענה על סדרות.
בעצם זיקקנו את ה\textquotedblright מהות\textquotedblright{} של המשפט
ושמנו אותה בשלב {\beginL 2\endL}, שהוא הפחות טכני; החלק הטכני הלך
ל-{\beginL 1\endL} שהוא בלתי נמנע כי אנחנו מדברים על מרחב \textbf{קונקרטי},
מרחב מטרי, ולכן משתמשים בטרמינולוגיה והכלים שלו. אבל גם שם, ההוכחה
מסובכת\textbf{ פחות} כי אם נסתכל על זה לרגע נראה שבהוכחה שמרחב מטרי
הוא האוסדורפי בגלל לא היינו צריכים לדבר על \textbf{סדרות}. סתם לקחנו
נקודה \L{$c$} כלשהי. אז מצד אחד פיצלנו את המשפט האחד לשני תתי-משפטים
שביחד הם יותר מסורבלים לתיאור מהמשפט המקורי; אבל כל חלק בנפרד הוא
יותר קל להבנה, וכל חלק בנפרד הוא בעל פוטנציאל לשימושים \textbf{נוספים.
}חלק {\beginL 1\endL} בעצם אומר לנו \textquotedblleft מעכשיו כל מה
שתוכיחו על מרחבים האוסדורפיים כלליים יחול אוטומטית על מרחבים מטריים\textquotedblright{}
וחלק {\beginL 2\endL} אומר לנו \textquotedblleft מעכשיו כל מרחב האוסדורפי
שתיתקלו בו יקיים את יחידות הגבול\textquotedblright .

זו הדגמה של מה שגורם לי להתלהב מטופולגיה קבוצתית - זה קורה גם בתחומים
מתמטיים אחרים כמובן, אבל כאן חזקה במיוחד התחושה של \textquotedblleft אוקיי
זה גם מגלה לי עולמות חדשים שלא חשבתי עליהם בכלל וגם גורם לי להסתכל
על מה שעשינו עד עכשיו באור חדש לגמרי\textquotedblright . והדוגמא הזו
היא רק קצה הקרחון.
\end{document}
